%! Graphing Polynomial Functions and Modeling — Unit 3, Lesson 3.6 — Warm-Up
\documentclass[10pt]{article}
\usepackage{algebra2-article}
\usepackage{algebra2-boxes}
\newcommand{\TallMath}[1]{$\displaystyle #1\rule[-1.4em]{0pt}{3.2em}$}

\begin{document}

\pageheader{Unit 3, Lesson 3.6}{Warm-Up}
\namedateperiod

\vspace{0.08in}

\begin{notesbox}{Getting Ready: Three Pieces of One Picture}
\small You already own every piece of a polynomial's graph. Today they get assembled. Each item
below is a tool you have used before --- notice what part of the picture it controls.

\vspace{5pt}
\textbf{1. The arms} (Lessons 3.0 and 3.1). You do \emph{not} have to multiply these out: degrees
add and leading coefficients multiply.

\vspace{3pt}
\renewcommand{\arraystretch}{1.5}
\begin{tabularx}{\linewidth}{|X|C{1.6cm}|C{2.0cm}|C{2.6cm}|C{2.6cm}|}
\hline
\rowcolor{forestbg}
\textbf{Function} & \textbf{Degree} & \textbf{Leading coeff.} & \textbf{As $x \to -\infty$, $f(x) \to$} & \textbf{As $x \to +\infty$, $f(x) \to$} \\ \hline
(a) $f(x) = (x-1)^2(x+2)$   & & & & \\ \hline
(b) $g(x) = -2x(x+3)^3$     & & & & \\ \hline
(c) $h(x) = -x^2(x-3)^3$    & & & & \\ \hline
\end{tabularx}

\vspace{6pt}
\textbf{2. The $x$-intercepts, and what happens at each one} (Lesson 3.0). Here is the polynomial
you have met in every lesson of this unit: \; $g(x) = x^3 - 3x + 2 = (x-1)^2(x+2)$.

\vspace{3pt}
Zeros: $x = $ \blank{1.2cm} with multiplicity \blank{0.9cm}, \quad and $x = $ \blank{1.2cm} with
multiplicity \blank{0.9cm}.

\vspace{3pt}
At a zero of \emph{even} multiplicity the graph \blank{2.4cm} the $x$-axis; at a zero of
\emph{odd} multiplicity it \blank{2.4cm} it.

\vspace{3pt}
The $y$-intercept: $g(0) = $ \blank{1.2cm}, so the graph passes through the point \blank{1.8cm}.

\vspace{6pt}
\textbf{3. Above or below?} The two zeros cut the $x$-axis into three stretches. Test one $x$-value
from each and record only the \emph{sign} of the answer.

\vspace{3pt}
\renewcommand{\arraystretch}{1.5}
\begin{tabularx}{\linewidth}{|C{3.2cm}|C{2.4cm}|C{2.6cm}|X|}
\hline
\rowcolor{forestbg}
\textbf{Interval} & \textbf{Test $x$} & \textbf{$g(x) = $} & \textbf{Sign of $g(x)$ --- is the graph above or below the $x$-axis?} \\ \hline
$x < -2$        & $-3$ & & \\ \hline
$-2 < x < 1$    & $0$  & & \\ \hline
$x > 1$         & $2$  & & \\ \hline
\end{tabularx}

\vspace{4pt}
\textbf{The hinge.} You now know which way both arms point, where the graph meets the $x$-axis and
whether it crosses or touches there, and whether it sits above or below the axis on every stretch in
between. \textbf{Could you pick this graph out of a lineup?} \blank{1.4cm}

\vspace{2pt}
What is the one thing this information still does \emph{not} tell you? \blank{6.2cm}

\end{notesbox}

\end{document}
